\PassOptionsToPackage{dvipsnames}{xcolor}
\documentclass[hyperref,UTF8,notheorems,xcolor={dvipsnames}]{beamer}
\usepackage[dvipsnames]{xcolor}
\usetheme{Rochester}
\usepackage{amsmath, mathtools, graphicx}
\usepackage{standalone}
\usepackage[most]{tcolorbox}
\usepackage[normalem]{ulem}
\usepackage{xeCJK}
\usepackage{fontspec}
\usepackage{tikz}
\usepackage{pgfplots}
\usepackage{pifont}
\usepackage{fontawesome5}
\usepgfplotslibrary{polar}
\usepgflibrary{shapes.geometric}
\usetikzlibrary{calc,angles,quotes}
\defaultfontfeatures{Mapping=tex-text}
\usefonttheme{professionalfonts}
\usepackage{concmath}
\usepackage{minted}

\newcommand{\cmark}{\ding{51}}%
\newcommand{\xmark}{\ding{55}}%

\newcommand{\CC}[1]{#1}

\newcommand{\Cincluded}{{\small\cmark}}
\newcommand{\Cdefine}{{\small\cmark\faFile*[regular]}}
\newcommand{\Ccodeonly}{{\small\cmark\faFile*~}}
\newcommand{\Cnofocus}{{\small\faQuestion}}
\newcommand{\Cexmaybe}{{\small\xmark\faQuestionCircle}}
\newcommand{\Cexcluded}{{\small\xmark}}

\newcommand{\Iincluded}{\item[\hbox to 1.8em{\Cincluded\hfill}]}
\newcommand{\Idefine}{\item[\hbox to 1.8em{\Cdefine\hfill}]}
\newcommand{\Icodeonly}{\item[\hbox to 1.8em{\Ccodeonly\hfill}]}
\newcommand{\Inofocus}{\item[\hbox to 1.8em{\Cnofocus\hfill}]}
\newcommand{\Iexmaybe}{\item[\hbox to 1.8em{\Cexmaybe\hfill}]}
\newcommand{\Iexcluded}{\item[\hbox to 1.8em{\Cexcluded\hfill}]}

\DeclareRobustCommand{\rddots}{\text{\reflectbox{$\ddots$}}}


\usepackage{graphicx}
\graphicspath{ {./../images/} }

\def\codesize{\fontsize{8}{9}\selectfont}
\setmonofont[Mapping=]{Noto Sans Mono}
\setminted{fontsize=\codesize, linenos, frame=lines, mathescape, autogobble, tabsize=4}
\setCJKmainfont[AutoFakeSlant,BoldFont=Noto Sans CJK TC Bold]{Noto Sans CJK TC}

\setlength{\parskip}{\baselineskip} 
\newcommand{\btitle}[1]{{\secname} -- #1}

\theoremstyle{definition}
\newtheorem{theorem}{定理}
\newtheorem{lemma}{引理}
\newtheorem{property}{性質}
\newtheorem{corollary}{推論}
\newtheorem{problem}{例題}


\newtheorem{definition}{定義}
\AtBeginEnvironment{definition}{
  \setbeamercolor{block title}{fg=white,bg=red!70!black}
  \setbeamercolor{block body}{fg=black, bg=block title.bg!10!bg}
}

\newtheorem{exercise}{習題}
\AtBeginEnvironment{exercise}{
  \setbeamercolor{block title}{fg=white,bg=green!30!black}
  \setbeamercolor{block body}{fg=black, bg=block title.bg!10!bg}
}

\newcommand{\sectionsubtitle}{} 
\newcommand{\setsubtitle}[1]{\renewcommand{\sectionsubtitle}{#1}}

\AtBeginSection[]{
  \begin{frame}
  \vfill
  \centering
  \begin{beamercolorbox}[sep=6pt,center,shadow=true,rounded=true]{title}
    \usebeamerfont{title}\LARGE\insertsectionhead\par%
    % 判斷 \sectionsubtitle 是否為空，如果不為空才渲染小字
    \ifx\sectionsubtitle\empty\else
      \vspace{4pt}% 微調大字與小字間距
      {\footnotesize\sectionsubtitle}\par%
    \fi
  \end{beamercolorbox}
  \vfill
  \end{frame}
  % 每次進入新 Section 後自動重置小字，避免影響下一個 Section
  \gdef\sectionsubtitle{}
}

\AtBeginSubsection[]{
  \begin{frame}
    \tableofcontents[subsectionstyle=show/shaded/hide]
  \end{frame}
}

\usepackage{ctable}
\usepackage{tabularx}

\setlength{\parskip}{\baselineskip}

\hypersetup{CJKbookmarks=true}

\usepackage{listings}
\usepackage{xcolor} % Required for defining custom colors

% Define color palette
\definecolor{codegreen}{rgb}{0,0.6,0}
\definecolor{codegray}{rgb}{0.5,0.5,0.5}
\definecolor{codepurple}{rgb}{0.58,0,0.82}
\definecolor{backcolour}{rgb}{0.95,0.95,0.92}

% Setup custom style block
\lstdefinestyle{mystyle}{
    backgroundcolor=\color{backcolour},
    commentstyle=\color{codegreen},
    keywordstyle=\color{magenta},
    numberstyle=\tiny\color{codegray},
    stringstyle=\color{codepurple},
    basicstyle=\ttfamily\footnotesize,
    breakatwhitespace=false,
    breaklines=true,                 % Auto-wrap long lines
    captionpos=b,                    % Put caption at bottom (b) or top (t)
    keepspaces=true,
    numbers=left,                    % Line numbers on the left
    numbersep=5pt,
    showspaces=false,
    showstringspaces=false,
    showtabs=false,
    tabsize=2
}

\lstset{style=mystyle} % Apply the style globally

